\chapter{Resultados y Discusión}
\label{cap:resultados_discusion}
En este capítulo se presentan los resultados obtenidos tras la ejecución de las fases experimentales descritas en la metodología. Se comienza con la validación del entorno de simulación para asegurar la fidelidad física de los modelos, seguido de la caracterización de los efectos degradantes en el $^{86}\text{Y}$ y, finalmente, la evaluación de las estrategias de optimización y la prueba de concepto de imagen multiplexada.

\section{Validación del Entorno de Simulación}
\label{sec:resultados_discusion:validacion}

Antes de abordar la optimización de parámetros, es imprescindible verificar que el simulador \texttt{penRed} y el módulo \texttt{PENNUC} reproducen correctamente la física de desintegración del isótopo y la geometría del escáner.

\subsection{Análisis del espectro de emisión del \yochoseis}
\label{subsec:resultados_espectro}

Para validar la fuente radiactiva, se extrajo el histograma de energía de los eventos individuales (\textit{singles}) detectados en una simulación sin medio dispersor (materiales \texttt{void}). La Figura \ref{fig:espectro_y86} muestra el espectro de energía obtenido.

\begin{figure}[htbp]
    \centering
    % Inserta aquí tu gráfica del espectro (la que tenía picos morados nítidos)
    \includegraphics[width=0.8\textwidth]{espectroy86.png}
    \caption{Espectro de energía de emisión del $^{86}\text{Y}$ obtenido mediante simulación Monte Carlo. Se observa el pico de aniquilación (\SI{511}{\kilo\electronvolt}) y las múltiples líneas de \textit{prompt gammas}.}
    \label{fig:espectro_y86}
\end{figure}

El espectro revela una serie de picos discretos y bien definidos. La ausencia de un fondo continuo Compton confirma que la simulación se realizó en vacío, registrando la emisión pura de la fuente. Se identifican claramente las emisiones características del $^{86}\text{Y}$:

\begin{enumerate}
    \item \textbf{Pico de Aniquilación (\SI{511}{\kilo\electronvolt}):} Correspondiente a la detección de los fotones generados por los positrones.
    \item \textbf{Picos de \textit{Prompt Gammas}:} Se observan emisiones de alta energía que coinciden con las transiciones nucleares del $^{86}\text{Sr}$.
\end{enumerate}

La Tabla \ref{tab:validacion_espectro} compara las energías de los picos observados en la simulación con los valores teóricos de la base de datos de referencia ENSDF (Evaluated Nuclear Structure Data File) consultada a través de NuDat 3 \parencite{NuDat3}.

\begin{table}[htbp]
    \centering
    \caption{Comparación entre las energías de emisión simuladas y los datos nucleares de referencia para el $^{86}\text{Y}$.}
    \label{tab:validacion_espectro}
    \begin{tabular}{c c c c}
    \hline
    \textbf{Origen} & \textbf{Energía Teórica (keV)} & \textbf{Energía Simulada (keV)} & \textbf{Intensidad Relativa (\%)} \\
    \hline
    Aniquilación $\beta^+$ & 511.0 & 511 & - \\
    Prompt $\gamma_1$ & 627.7 & $\sim 628$ & 32.6 \\
    Prompt $\gamma_2$ & 1076.6 & $\sim 1077$ & 82.5 \\
    Prompt $\gamma_3$ & 1153.1 & $\sim 1153$ & 30.5 \\
    Prompt $\gamma_4$ & 1854.4 & $\sim 1854$ & 17.2 \\
    Prompt $\gamma_5$ & 1920.7 & $\sim 1921$ & 20.8 \\
    \hline
    \end{tabular}
\end{table}

La coincidencia precisa entre los picos simulados y los teóricos valida el uso del archivo \texttt{Y86.nuc} y el correcto funcionamiento del módulo \texttt{PENNUC} para generar las cascadas gamma que causan las coincidencias triples. Es notable la alta intensidad del pico de \SI{1077}{\kilo\electronvolt}, que supera en probabilidad de emisión por desintegración a los propios fotones de aniquilación, anticipando su gran impacto en la señal de fondo.

\subsection{Resultados de la simulación ideal (\textit{Ground Truth})}
\label{subsec:resultados_ideal}

En la Fase 1 del diseño experimental, se realizó la reconstrucción de la imagen en condiciones ideales: materiales transparentes (vacío) y filtrado de coincidencias mediante el historial de la partícula (modo \textit{Same History}). El resultado se muestra en la Figura \ref{fig:imagen_ideal}.

\begin{figure}[htbp]
    \centering
    % Inserta aquí la imagen del fantoma perfecto (la que tenía fondo negro puro)
    \includegraphics[width=0.6\textwidth]{IMG_IDEAL_Y86.png}
    \caption{Imagen reconstruida de la simulación ideal (\textit{Ground Truth}). Se observa una definición geométrica perfecta de las varillas del fantoma Derenzo y una ausencia total de ruido de fondo.}
    \label{fig:imagen_ideal}
\end{figure}

El análisis visual y cuantitativo de esta imagen establece la referencia de máxima calidad alcanzable:

\begin{itemize}
    \item \textbf{Geometría:} La distribución espacial de la actividad reproduce fielmente la geometría del fantoma Derenzo modelado. Se distinguen claramente los seis sectores de varillas, incluyendo las de menor diámetro (\SI{0.8}{\milli\meter}), confirmando la correcta definición del voxelizado y la proyección en el software de reconstrucción.
    
    \item \textbf{Ausencia de Ruido:} El fondo de la imagen presenta una intensidad nula o despreciable. Esto valida el algoritmo de post-procesado en modo \textit{Same History}, demostrando que es capaz de filtrar el 100\% de las coincidencias aleatorias y los eventos espurios generados por los \textit{prompt gammas} no correlacionados.
    
    \item \textbf{Métricas de Referencia:} El análisis de ROIs arrojó un \textbf{Coeficiente de Recuperación de Contraste (CRC) de 185.27}. Este valor extremadamente alto (técnicamente infinito si el fondo fuera cero absoluto) sirve como el estándar de oro (\textit{Gold Standard}) contra el cual se comparará la degradación en las fases realistas.
\end{itemize}

Estos resultados confirman que la cadena completa de simulación —desde la generación de partículas en \texttt{penRed} hasta la reconstrucción en \texttt{RecoServer}— funciona correctamente y que cualquier degradación observada en fases posteriores será consecuencia directa de los fenómenos físicos introducidos y no de errores de software.
\section{Caracterización de la Degradación de la Imagen}
\label{sec:resultados_discusion:scater_randoms}

Una vez establecida la referencia ideal, las siguientes fases del estudio introdujeron progresivamente los efectos físicos que ocurren en una adquisición real. El objetivo es aislar y cuantificar la contribución de cada fenómeno a la degradación de la calidad de la imagen.

\subsection{Impacto aislado del Scatter y la Atenuación}
\label{subsec:resultados_scatter}

En la Fase 2, se activó la física de interacción con la materia al definir el fantoma como agua, pero se mantuvo el filtro ideal de coincidencias (\textit{Take Same History}). La Figura \ref{fig:imagen_scatter} compara la imagen resultante con la de la fase ideal.

\begin{figure}[htbp]
    \centering
    \begin{subfigure}[b]{0.48\textwidth}
        \centering
        \includegraphics[width=\textwidth]{IMG_IDEAL_Y86.png}
        \caption{Imagen Ideal (Fase 1).}
        \label{fig:scatter_vs_ideal_a}
    \end{subfigure}
    \hfill
    \begin{subfigure}[b]{0.48\textwidth}
        \centering
        \includegraphics[width=\textwidth]{IMG_SCATTER_Y86.png}
        \caption{Imagen con Scatter y Atenuación (Fase 2).}
        \label{fig:scatter_vs_ideal_b}
    \end{subfigure}
    \caption{Comparación visual del impacto del \textit{scatter} y la atenuación. La imagen de la derecha muestra una pérdida de resolución (halos) y un fondo elevado que reduce el contraste.}
    \label{fig:imagen_scatter}
\end{figure}

El análisis revela dos efectos principales:
\begin{itemize}
    \item \textbf{Dispersión Compton (\textit{Scatter}):} La interacción de los fotones con el agua provoca una dispersión que genera \enquote{halos} alrededor de las varillas calientes y un fondo difuso en toda la imagen. Esto se traduce en una pérdida de resolución espacial y una reducción del contraste.
    \item \textbf{Atenuación:} Una fracción significativa de los fotones es absorbida o dispersada fuera del campo de visión, lo que reduce el número total de cuentas verdaderas y, por tanto, incrementa el ruido estadístico.
\end{itemize}

Cuantitativamente, el CRC se desplomó de \textbf{185.27} en la fase ideal a \textbf{11.57} en esta fase. Esto representa una pérdida de contraste del 93.7\% atribuible exclusivamente a la física de interacción en el objeto y a los \textit{prompt gammas} no filtrados por el procesado ideal.

\subsection{Impacto de las Coincidencias Aleatorias}
\label{subsec:resultados_randoms}

En la Fase 3, se introdujo el comportamiento de un detector realista, desactivando el filtro por historial (\textit{Take Closest}). Esto añade las coincidencias aleatorias (\textit{randoms}) a los efectos de \textit{scatter} y atenuación ya presentes.

El resultado es la imagen de referencia realista (IMG\_CLOSEST\_TIME-Y86), que representa la calidad de imagen obtenida en un escáner estándar sin correcciones avanzadas. Visualmente, el fondo de la imagen se vuelve significativamente más ruidoso y elevado, degradando aún más el contraste.
Cuantitativamente, el CRC se mantuvo en un valor similar de \textbf{11.73}, mientras que el SNR empeoró ligeramente a \textbf{2.74}, confirmando que el mayor daño al contraste proviene de los eventos de \textit{scatter} y \textit{prompt gammas} correlacionados, y que los \textit{randoms} contribuyen principalmente al ruido estadístico general.

\subsection{Comparativa crítica: \yochoseis vs \fdiezocho}
\label{subsec:resultados_comparativa}

Para aislar el impacto específico de la física de desintegración del $^{86}\text{Y}$, se realizó una simulación de control en condiciones realistas (Fase 3) utilizando un isótopo \enquote{limpio}, el $^{18}\text{F}$. La Figura \ref{fig:comparativa_y86_f18} muestra la comparación directa entre ambas imágenes.

\begin{figure}[htbp]
    \centering
    \begin{subfigure}[b]{0.48\textwidth}
        \centering
        \includegraphics[width=\textwidth]{IMG_CLOSEST_TIME-F18.png}
        \caption{Imagen Realista con $^{18}\text{F}$.}
        \label{fig:comparativa_f18}
    \end{subfigure}
    \hfill
    \begin{subfigure}[b]{0.48\textwidth}
        \centering
        \includegraphics[width=\textwidth]{IMG_CLOSEST_TIME-Y86.png}
        \caption{Imagen Realista con $^{86}\text{Y}$.}
        \label{fig:comparativa_y86}
    \end{subfigure}
    \caption{Comparación de la calidad de imagen en condiciones realistas. La imagen de $^{18}\text{F}$ presenta un fondo significativamente más limpio y un contraste superior, demostrando el impacto negativo de los \textit{prompt gammas} del $^{86}\text{Y}$.}
    \label{fig:comparativa_y86_f18}
\end{figure}

El análisis cuantitativo de esta comparación es concluyente:

\begin{table}[htbp]
    \centering
    \caption{Comparación de métricas de calidad entre $^{86}\text{Y}$ y $^{18}\text{F}$ en condiciones realistas.}
    \label{tab:comparativa_metricas}
    \begin{tabular}{c c c}
    \hline
    \textbf{Isótopo} & \textbf{CRC} & \textbf{SNR} \\
    \hline
    $^{18}\text{F}$ & 31.09 & 3.86 \\
    $^{86}\text{Y}$ & 11.73 & 2.74 \\
    \hline
    \textbf{Diferencia} & \textbf{2.6x Menos Contraste} & \textbf{29\% Menos SNR} \\
    \hline
    \end{tabular}
\end{table}

\textbf{Conclusión de la caracterización:}
La dramática diferencia en calidad de imagen, con una reducción del contraste en un factor de 2.6, no puede explicarse únicamente por el \textit{scatter} o los \textit{randoms} convencionales, ya que ambos afectan a los dos isótopos.
Esta evidencia demuestra que el principal factor de degradación en la imagen de $^{86}\text{Y}$ es el fondo masivo de coincidencias espurias (aleatorias y correlacionadas) generado por la alta emisión de \textit{prompt gammas}, justificando la necesidad de las estrategias de optimización que se presentan a continuación.
Una vez caracterizada la degradación de la imagen en condiciones realistas, el siguiente paso fue implementar y evaluar estrategias para mitigarla. Esta sección presenta los resultados de la optimización de los parámetros de post-procesado (Ventanas de Energía y Tiempo) sobre la imagen de referencia realista del $^{86}\text{Y}$ (Fase 3), utilizando como base una simulación de $1.5 \cdot 10^9$ historias para asegurar una estadística robusta.

\subsection{Efecto de la Ventana de Energía (EW)}
\label{subsec:resultados_ew}

El principal mecanismo para rechazar eventos espurios es el filtrado energético. Se evaluaron tres ventanas progresivamente más estrechas, manteniendo una ventana de coincidencia fija de \SI{4}{\nano\second}. Los resultados, tanto visuales como cuantitativos, se muestran en la Figura \ref{fig:resultados_ew} y la Tabla \ref{tab:resultados_ew}.

\begin{figure}[htbp]
    \centering
    \begin{subfigure}[b]{0.32\textwidth}
        \centering
        \includegraphics[width=\textwidth]{IMG_OPT_EW1_CW1.png}
        \caption{EW 20\% (409-613 keV).}
    \end{subfigure}
    \hfill
    \begin{subfigure}[b]{0.32\textwidth}
        \centering
        \includegraphics[width=\textwidth]{IMG_OPT_EW2_CW1.png}
        \caption{EW 15\% (434-588 keV).}
    \end{subfigure}
    \hfill
    \begin{subfigure}[b]{0.32\textwidth}
        \centering
        \includegraphics[width=\textwidth]{IMG_OPT_EW3_CW1.png}
        \caption{EW 10\% (460-562 keV).}
    \end{subfigure}
    \caption{Impacto visual de la reducción de la Ventana de Energía (EW). A medida que la ventana se estrecha, se observa una drástica reducción del ruido de fondo y una clara recuperación de la visibilidad de las varillas más pequeñas.}
    \label{fig:resultados_ew}
\end{figure}

\begin{table}[htbp]
    \centering
    \caption{Métricas de calidad de imagen en función de la Ventana de Energía (CW = \SI{4}{\nano\second}).}
    \label{tab:resultados_ew}
    \begin{tabular}{c c c}
    \hline
    \textbf{Ventana de Energía} & \textbf{CRC} & \textbf{SNR} \\
    \hline
    30\% (Línea Base) & 11.73 & 2.74 \\
    20\% & 10.19 & 6.89 \\
    15\% & 13.06 & 6.32 \\
    10\% & 19.17 & 6.64 \\
    \hline
    \end{tabular}
\end{table}

El análisis de estos resultados revela que estrechar la ventana de energía es una estrategia altamente efectiva. Al pasar de una ventana del 30\% a una del 10\%, se logra un incremento del 63\% en el Contraste (CRC), superando incluso el valor de la línea base. Es notable que el SNR también mejora masivamente, indicando que el rechazo de eventos de fondo es tan significativo que compensa la pérdida de cuentas verdaderas.

\subsection{Efecto de la Ventana de Coincidencia (CW)}
\label{subsec:resultados_cw}

A continuación, se evaluó el impacto de acortar la Ventana de Coincidencia (CW) para reducir las coincidencias aleatorias. Se comparó una CW de \SI{4}{\nano\second} con una de \SI{3}{\nano\second}, manteniendo fija la ventana de energía en las tres configuraciones anteriores.

Los resultados (Tabla \ref{tab:resultados_cw_comparativa}) muestran que acortar la ventana temporal de 4 a 3 ns produce una mejora marginal pero consistente en el SNR (del orden del 1-5\%) debido a la reducción de \textit{randoms}, mientras que el impacto en el CRC es mínimo y no sigue una tendencia clara. Esto indica que la ventana de energía es el parámetro dominante para la recuperación del contraste.

\begin{table}[htbp]
    \centering
    \caption{Comparación de métricas al reducir la Ventana de Coincidencia (CW).}
    \label{tab:resultados_cw_comparativa}
    \begin{tabular}{c | c c | c c}
    \hline
    \textbf{Ventana de Energía} & \multicolumn{2}{c|}{\textbf{CW = 4 ns}} & \multicolumn{2}{c}{\textbf{CW = 3 ns}} \\
     & \textbf{CRC} & \textbf{SNR} & \textbf{CRC} & \textbf{SNR} \\
    \hline
    20\% & 10.19 & 6.89 & 9.99 & 6.87 \\
    15\% & 13.06 & 6.32 & 13.75 & 6.72 \\
    10\% & 19.17 & 6.64 & 18.04 & 6.72 \\
    \hline
    \end{tabular}
\end{table}

\subsection{Análisis de Trade-off y Determinación del Punto Óptimo}
\label{subsec:resultados_tradeoff}

Para determinar la mejor configuración global, se deben analizar conjuntamente el CRC y el SNR. La Figura \ref{fig:tradeoff} representa el \enquote{trade-off} entre el contraste recuperado y el ruido de la imagen para todas las configuraciones realistas.

% \begin{figure}[htbp]
%     \centering
%     % Inserta aquí tu gráfico de dispersión CRC vs SNR
%     \includegraphics[width=0.7\textwidth]{grafico_tradeoff.png}
%     \caption{Gráfico de compromiso (trade-off) CRC vs. SNR. Cada punto representa una configuración de optimización. El objetivo es alcanzar el punto más alto y a la derecha posible.}
     \label{fig:tradeoff}
% \end{figure}

El análisis del gráfico revela que la configuración \textbf{EW 10\% / CW 4ns} (\texttt{IMG\_OPT\_EW3\_CW1}) ofrece el mejor compromiso. Alcanza el máximo valor de CRC (19.17), fundamental para la detectabilidad de lesiones pequeñas, mientras mantiene un SNR excelente (6.64), muy superior a la línea base. Aunque otras configuraciones ofrecen un SNR marginalmente superior, lo hacen a costa de una pérdida significativa de contraste.

\subsection{Convergencia y Resultados de Alta Estadística}
\label{subsec:resultados_convergencia}

Finalmente, para validar la robustez de la configuración óptima y obtener una imagen de máxima calidad, se realizó una simulación final (\texttt{SIM\_FINAL}) con una estadística 4 veces superior ($6 \cdot 10^9$ historias). Se aplicaron los parámetros óptimos de la configuración \texttt{IMG\_OPT\_EW2\_CW1} (EW 15\% / CW 4ns).

Los resultados de esta simulación de alta estadística (CRC = 8.13, SNR = 7.91) muestran una mejora del 25\% en el SNR en comparación con su homóloga de baja estadística (SNR = 6.32), confirmando que el ruido estadístico se reduce con el aumento de cuentas.

La aparente disminución del CRC (de 13.06 a 8.13) no debe interpretarse como un empeoramiento, sino como una convergencia hacia el valor de contraste real del sistema. En simulaciones con menos estadística, el ruido puede inflar o desinflar artificialmente el valor del CRC. La simulación de alta estadística proporciona una medida más precisa y fiable, estableciendo el verdadero rendimiento cuantitativo de la configuración optimizada.

\section{Implementación de Imagen Multiplexada}
\label{sec:resultados_discusion:multiplexing}

Una vez optimizada la adquisición para un solo isótopo, la última fase de este trabajo explora el potencial del $^{86}\text{Y}$ como habilitador de la imagen PET multiplexada (mPET). El objetivo de esta prueba de concepto es demostrar la capacidad de separar, a partir de una única adquisición simultánea, las señales de un isótopo \enquote{sucio} ($^{86}\text{Y}$) y uno \enquote{limpio} ($^{18}\text{F}$).

Se simuló un fantoma con dos cilindros activos: el de la izquierda contenía exclusivamente $^{18}\text{F}$ y el de la derecha, $^{86}\text{Y}$.

\subsection{Identificación y Separación de Eventos Triples}
\label{subsec:resultados_triples}

La clave de la separación reside en el algoritmo de post-procesado desarrollado en este TFG (ver sección \ref{sec:metodologia:postprocess}), ejecutado en modo \textit{Multiplexing}. Este algoritmo clasifica cada coincidencia detectada basándose en si existe o no un tercer fotón simultáneo con alta energía (un \textit{prompt gamma}).

El resultado de este proceso es la generación de dos \textit{datasets} en modo lista, completamente independientes:
\begin{enumerate}
    \item \textbf{Dataset de Dobles (A+B):} Contiene todas las coincidencias de dos fotones. Incluye los eventos verdaderos del $^{18}\text{F}$ y las coincidencias dobles del $^{86}\text{Y}$ (tanto verdaderas como de \textit{scatter} y \textit{randoms}).
    \item \textbf{Dataset de Triples (B'):} Contiene exclusivamente los eventos donde se detectó el tercer fotón. Como el $^{18}\text{F}$ no emite \textit{prompt gammas}, este dataset representa una señal pura del $^{86}\text{Y}$.
\end{enumerate}

\subsection{Recuperación de las Imágenes de Isótopos Individuales}
\label{subsec:resultados_recuperacion}

Las imágenes reconstruidas a partir de estos dos \textit{datasets} se muestran en la Figura \ref{fig:multiplexing_inicial}. La imagen de dobles (A+B) muestra la actividad mezclada, mientras que la imagen de triples (B') aísla con éxito la señal del $^{86}\text{Y}$, aunque con una estadística considerablemente menor.

\begin{figure}[htbp]
    \centering
    \includegraphics[width=0.9\textwidth]{imagen_cuatro_subplots.png} % Tu imagen con los 4 subplots
    \caption{Resultados de la separación de isótopos. \textbf{AB:} Imagen de dobles mixta. \textbf{B':} Imagen de triples, que aísla el $^{86}\text{Y}$. \textbf{B:} Estimación de la contaminación de dobles del $^{86}\text{Y}$. \textbf{A:} Imagen recuperada del $^{18}\text{F}$ tras la sustracción.}
    \label{fig:multiplexing_inicial}
\end{figure}

Para recuperar la imagen pura del $^{18}\text{F}$, se aplicó un algoritmo de sustracción lineal en Python. Primero, se estimó la contribución de dobles del $^{86}\text{Y}$ (Imagen B) escalando la imagen de triples (B') con un factor $\beta$ determinado experimentalmente. El factor $\beta$, que representa el ratio de dobles a triples para el $^{86}\text{Y}$ en este sistema, se calculó en 75.8.

\begin{equation}
    I_{Y86}^{\text{D}} = \beta \cdot I_{Y86}^{\text{T}}
\end{equation}

Finalmente, esta contribución estimada se sustrajo de la imagen mixta total:

\begin{equation}
    I_{F18} = I_{AB}^{\text{D}} - I_{Y86}^{\text{D}}
\end{equation}

\subsection{Resultados Visuales de la Separación}
\label{subsec:resultados_visuales_separacion}

El resultado final del proceso se muestra en la Figura \ref{fig:multiplexing_final}. La imagen recuperada del $^{18}\text{F}$ (Imagen A) muestra actividad únicamente en el cilindro de la izquierda, mientras que la señal del cilindro del $^{86}\text{Y}$ ha sido eliminada casi por completo.

\begin{figure}[htbp]
    \centering
    \begin{subfigure}[b]{0.48\textwidth}
        \centering
        \includegraphics[width=0.8\textwidth]{subplot_A.png} % Tu subplot de la imagen A recuperada
        \caption{Imagen Recuperada del $^{18}\text{F}$ (A).}
    \end{subfigure}
    \hfill
    \begin{subfigure}[b]{0.48\textwidth}
        \centering
        \includegraphics[width=0.8\textwidth]{subplot_B_prime.png} % Tu subplot de la imagen B_prime
        \caption{Imagen Aislada del $^{86}\text{Y}$ (B').}
    \end{subfigure}
    \caption{Imágenes finales separadas. El algoritmo ha logrado aislar la distribución espacial de cada isótopo a partir de una única adquisición simultánea.}
    \label{fig:multiplexing_final}
\end{figure}


Estos resultados constituyen una \textbf{prueba de concepto exitosa}, demostrando que el método de clasificación de eventos triples es una estrategia viable para la implementación de la imagen PET multiplexada. Se valida así el potencial de los isótopos \enquote{sucios} no como una fuente de problemas, sino como una herramienta para expandir las capacidades de la imagen PET.